\section{여러 원소로 생성한 확장과 대수성의 추이성}
\label{sec:generated-algebraic-extensions}

\begin{learninggoal}
유한 개의 대수적 원소가 생성하는 확장이 유한확장임을 증명하고, 대수적 확장의 추이성을 이해한다. 유한확장과 대수적 확장의 차이를 구별한다.
\end{learninggoal}

\subsection{여러 원소를 차례로 첨가하기}

체확장 $L/K$와 원소 $\alpha_1,\dots,\alpha_r\in L$에 대해
\[
  K(\alpha_1,\dots,\alpha_r)
\]
는 $K$와 이 원소들을 포함하는 가장 작은 부분체이다. 이를 한 번에 만드는 대신
\[
  K\subseteq K(\alpha_1)\subseteq K(\alpha_1,\alpha_2)
  \subseteq\cdots\subseteq K(\alpha_1,\dots,\alpha_r)
\]
와 같이 하나씩 첨가할 수 있다.

\begin{theorem}
$\alpha_1,\dots,\alpha_r$가 모두 $K$ 위에서 대수적이면
\[
  K(\alpha_1,\dots,\alpha_r)/K
\]
는 유한확장이다. 더 정확히,
\[
  [K(\alpha_1,\dots,\alpha_r):K]
  \le\prod_{j=1}^r\deg m_{\alpha_j,K}.
\]
\end{theorem}

\begin{proof}
각 $j$에 대해
\[
  K_{j-1}=K(\alpha_1,\dots,\alpha_{j-1}),
  \qquad
  K_j=K_{j-1}(\alpha_j)
\]
라 하자. $\alpha_j$는 $K$ 위에서 대수적이므로 같은 소거다항식을 $K_{j-1}[X]$에서도 만족한다. 따라서 $\alpha_j$는 $K_{j-1}$ 위에서 대수적이고
\[
  [K_j:K_{j-1}]\le\deg m_{\alpha_j,K}.
\]
탑 법칙을 반복 적용하면
\[
  [K_r:K]
  =\prod_{j=1}^r[K_j:K_{j-1}]
  \le\prod_{j=1}^r\deg m_{\alpha_j,K}.
\]
\end{proof}

\begin{corollary}
유한확장 $L/K$에 대해 다음은 동치이다.
\begin{enumerate}[label=(\roman*)]
  \item $L/K$는 유한확장이다.
  \item $L=K(\alpha_1,\dots,\alpha_r)$이고 각 $\alpha_j$가 $K$ 위에서 대수적인 유한 생성계가 존재한다.
  \item $L/K$는 유한생성 대수적 확장이다.
\end{enumerate}
\end{corollary}

\begin{proof}
(i)이면 $K$-기저를 생성계로 택할 수 있고, 유한확장은 대수적이다. (ii)에서 (iii)은 정의이며, (iii)에서 (i)는 앞의 정리이다.
\end{proof}

\subsection{예제: 두 제곱근을 첨가하기}

\begin{proposition}
\[
  [\Q(\sqrt2,\sqrt3):\Q]=4
\]
이고
\[
  1,\sqrt2,\sqrt3,\sqrt6
\]
이 $\Q$-기저이다.
\end{proposition}

\begin{proof}
먼저 $[\Q(\sqrt2):\Q]=2$이다. 이제 $\sqrt3\notin\Q(\sqrt2)$임을 보이자. 만약
\[
  \sqrt3=a+b\sqrt2
  \qquad(a,b\in\Q)
\]
이면 제곱하여
\[
  3=a^2+2b^2+2ab\sqrt2
\]
를 얻는다. $1,\sqrt2$의 선형독립성에서 $ab=0$이다. $b=0$이면 $\sqrt3\in\Q$가 되고, $a=0$이면 $b^2=3/2$인데 유리수 제곱의 소인수 지수는 모두 짝수이므로 불가능하다. 따라서 $\sqrt3$의 $\Q(\sqrt2)$ 위 최소다항식은 $X^2-3$이고
\[
  [\Q(\sqrt2,\sqrt3):\Q(\sqrt2)]=2.
\]
탑 법칙으로 전체 차수는 $4$이다. 두 단계의 기저를 곱하면 제시한 네 원소가 기저가 된다.
\end{proof}

\subsection{대수성은 확장탑을 따라 전해진다}

\begin{theorem}[대수성의 추이성]
체확장탑
\[
  K\subseteq L\subseteq M
\]
에서 $L/K$와 $M/L$이 대수적이면 $M/K$도 대수적이다.
\end{theorem}

\begin{proof}
$\alpha\in M$를 택하자. $\alpha$는 $L$ 위에서 대수적이므로 어떤 일계수다항식
\[
  X^n+c_{n-1}X^{n-1}+\cdots+c_0\in L[X]
\]
의 근이다. 계수 $c_0,\dots,c_{n-1}$은 $L/K$가 대수적이므로 모두 $K$ 위에서 대수적이다. 따라서
\[
  E=K(c_0,\dots,c_{n-1})
\]
는 $K$의 유한확장이다. 한편 $\alpha$는 $E$ 위에서 대수적이므로 $E(\alpha)/E$도 유한하다. 탑 법칙에서 $E(\alpha)/K$가 유한하고, 유한확장은 대수적이므로 $\alpha$는 $K$ 위에서 대수적이다.
\end{proof}

\begin{corollary}
확장 $L/K$에서 $K$ 위 대수적인 원소들의 집합
\[
  A=\{\alpha\in L:\alpha\text{는 }K\text{ 위에서 대수적}\}
\]
은 $L$의 부분체이다.
\end{corollary}

\begin{proof}
$\alpha,\beta\in A$이면 $K(\alpha,\beta)/K$가 유한확장이다. 따라서 그 안의
\[
  \alpha\pm\beta,
  \qquad
  \alpha\beta,
  \qquad
  \alpha/\beta\quad(\beta\ne0)
\]
도 $K$ 위에서 대수적이다.
\end{proof}

\subsection{대수적이지만 유한하지 않은 확장}

유한확장은 대수적이지만 그 역은 거짓이다. 복소수 안의 대수적 수 전체를
\[
  \overline\Q_{\mathrm{alg}}
  =\{z\in\C:z\text{는 }\Q\text{ 위에서 대수적}\}
\]
라 하자. 앞의 따름정리에 의해 이는 체이며 $\overline\Q_{\mathrm{alg}}/\Q$는 대수적 확장이다.

그러나 이 확장은 유한하지 않다. 각 양의 정수 $n$에 대해 아이젠슈타인 판정법으로
\[
  X^n-2
\]
는 $\Q[X]$에서 기약이므로
\[
  [\Q(\sqrt[n]{2}):\Q]=n.
\]
만약 $[\overline\Q_{\mathrm{alg}}:\Q]$가 유한이라면 모든 중간체의 차수는 그 유한한 수를 넘을 수 없으므로 모순이다.

\begin{pitfall}
``대수적''은 각 원소가 어떤 다항식 관계를 만족한다는 조건이고, ``유한''은 전체 확장의 벡터공간 차원이 유한하다는 조건이다. 원소별로는 모두 대수적이어도 필요한 차수들이 무한히 커질 수 있다.
\end{pitfall}

\begin{checkpoint}
\begin{enumerate}[label=(\alph*)]
  \item $[\Q(\sqrt2,\sqrt5):\Q]$를 계산하고 기저를 제시하라.
  \item $\alpha,\beta$가 $K$ 위에서 대수적이면 $\alpha+\beta$도 대수적임을 유한확장으로 설명하라.
  \item 대수적이지만 유한하지 않은 확장의 예를 하나 설명하라.
\end{enumerate}
\end{checkpoint}
